% =====================================================================
%  preamble.tex — 压缩与验证：复杂软件系统的层级涌现、结构收敛与可验证执行
%  编译方式：xelatex main.tex （连续两次，以生成目录与交叉引用）
%  依赖：TeX Live 2019 及以上完整发行版（ctex / amsmath / amsthm / booktabs /
%        longtable / enumitem / tikz / hyperref）
% =====================================================================

% ---------- 页面 ----------
\usepackage{geometry}
\geometry{a4paper,left=2.7cm,right=2.7cm,top=2.9cm,bottom=2.9cm}

% 减少中文长行导致的 overfull box
\setlength{\emergencystretch}{2em}
\tolerance=2000
\hbadness=10000

% ---------- 数学 ----------
\usepackage{amsmath}
\usepackage{amssymb}
\usepackage{mathtools}
\usepackage{amsthm}

% ---------- 表格与列表 ----------
\usepackage{array}
\setlength{\tabcolsep}{4pt}
\usepackage{booktabs}
\usepackage{longtable}
\usepackage{enumitem}
\setlist{nosep,leftmargin=2.2em}

% ---------- 图形 ----------
\usepackage{tikz}
\usetikzlibrary{arrows.meta,positioning,fit,backgrounds,calc}

% ---------- 颜色与标题 ----------
\usepackage{xcolor}
\usepackage{caption}
\captionsetup{font=small,labelfont=bf}

% ---------- 定理环境 ----------
% 行内小标题样式：用于承载论断性而非纯数学性的条目，
% 使其以论文段首小标题的形式排版，而非独立展示块。
\newtheoremstyle{runin}%
  {0.7em}%      上间距
  {0.2em}%      下间距
  {\normalfont}% 正文
  {}%           缩进
  {\bfseries}%  标题字体
  {.}%          标题后标点
  {0.5em}%      标题后间距
  {}%           标题格式（默认）

\theoremstyle{plain}
\newtheorem{theorem}{定理}[chapter]
\newtheorem{proposition}[theorem]{命题}
\newtheorem{lemma}[theorem]{引理}
\newtheorem{corollary}[theorem]{推论}
\newtheorem{conjecture}[theorem]{猜想}

\theoremstyle{definition}
\newtheorem{definition}[theorem]{定义}
\newtheorem{hypothesis}[theorem]{假设}
\newtheorem{example}[theorem]{例}
\newtheorem{problem}[theorem]{问题}

\theoremstyle{runin}
\newtheorem{principle}[theorem]{原则}

\theoremstyle{remark}
\newtheorem{remark}[theorem]{注}
\newtheorem{observation}[theorem]{观察}

\renewcommand{\proofname}{证明}
\renewcommand{\qedsymbol}{$\blacksquare$}

% ---------- 常用命令 ----------
\newcommand{\Beh}{\mathcal{B}}
\newcommand{\Obs}{\mathcal{O}}
\newcommand{\Env}{\mathcal{A}}
\newcommand{\Cost}{\mathcal{C}}
\newcommand{\Inv}{\mathcal{I}}
\newcommand{\Mod}{\mathcal{M}}
\newcommand{\Sem}{\llbracket}
\newcommand{\Ese}{\rrbracket}
\newcommand{\Tr}{\operatorname{Traces}}
\newcommand{\Allow}{\operatorname{Allowed}}
\newcommand{\Observe}{\alpha}
\newcommand{\Conc}{\gamma}
\newcommand{\Obl}{\operatorname{Obl}}
\newcommand{\Ker}{\operatorname{ker}}
\newcommand{\TCB}{\ensuremath{\mathrm{TCB}}}
\newcommand{\DL}{\ell}
\newcommand{\Entropy}{H}
\newcommand{\BigO}{\mathcal{O}}
\newcommand{\code}[1]{\texttt{#1}}
\newcommand{\term}[1]{\textbf{#1}}
\newcommand{\cn}[1]{\textnormal{#1}}

% 层级宏：\Lk{0} 生成 $\mathrm{L}_0$
\newcommand{\Lk}[1]{\ensuremath{\mathrm{L}_{#1}}}

% 便于排版的列类型
\newcolumntype{L}[1]{>{\raggedright\arraybackslash}p{#1}}
\newcolumntype{C}[1]{>{\centering\arraybackslash}p{#1}}

% ---------- 超链接（最后加载） ----------
\usepackage[colorlinks=true,
            linkcolor=blue!55!black,
            citecolor=teal!70!black,
            urlcolor=blue!65!black,
            bookmarksnumbered=true,
            unicode=true]{hyperref}
